If $xa=0$ with $a\ne0$, choose a prime minimal over $(0:a)$; it contains $x$ and belongs to $D(A)$. Conversely, suppose $x\in\mathfrak{p}$ and $\mathfrak{p}$ is minimal over $(0:a)$. In $A_{\mathfrak{p}}$, the radical of $(0:a)A_{\mathfrak{p}}$ is $\mathfrak{p}A_{\mathfrak{p}}$, so $tx^na=0$ for some $t\notin\mathfrak{p}$ and $n>0$. Since $ta\ne0$, choosing the least such exponent shows that $x$ annihilates a nonzero element.

We have $\operatorname{Ann}_{S^{-1}A}(a/s)=S^{-1}(0:a)$. A prime minimal over $(0:a)$ and disjoint from $S$ extends to a prime minimal over this annihilator. Conversely, if a prime of $S^{-1}A$ is minimal over this annihilator, choose a prime minimal over $(0:a)$ inside its contraction. This prime is disjoint from $S$, and minimality after extension shows that it equals the contraction. This proves the asserted equality.

If $(0)=\bigcap_{i=1}^n\mathfrak{q}_i$ is a minimal primary decomposition, with $\mathfrak{p}_i=r(\mathfrak{q}_i)$, the first uniqueness theorem gives $\mathfrak{p}_i=r(0:a)$ for a suitable $a$, so $\mathfrak{p}_i\in D(A)$. Conversely,
\[r(0:a)=\bigcap_{a\notin\mathfrak{q}_i}\mathfrak{p}_i.\]
A prime minimal over $(0:a)$ contains one of the primes on the right, which also contains $(0:a)$. Minimality makes them equal.
